\documentclass[11pt,a4paper]{article}

% ================ Encodage et langue ================
\usepackage[utf8]{inputenc}
\usepackage[T1]{fontenc}
\usepackage[french]{babel}
\usepackage{lmodern}
\usepackage{textcomp}

% ================ Mise en page ================
\usepackage[margin=2.2cm]{geometry}
\usepackage{parskip}
\setlength{\parindent}{0pt}
\setlength{\parskip}{0.6em}

% ================ Couleurs et liens ================
\usepackage[table]{xcolor}
\definecolor{primary}{HTML}{0B5394}
\definecolor{accent}{HTML}{1F77B4}
\definecolor{light}{HTML}{F2F2F2}
\definecolor{codebg}{HTML}{F6F8FA}
\definecolor{rulec}{HTML}{D0D7DE}

\usepackage[colorlinks=true, linkcolor=primary, urlcolor=primary, citecolor=primary]{hyperref}

% ================ Tableaux ================
\usepackage{longtable, booktabs, array, tabularx}
\renewcommand{\arraystretch}{1.25}
\arrayrulecolor{rulec}

% ================ Listes ================
\usepackage{enumitem}
\setlist{itemsep=2pt, topsep=4pt}

% ================ Code et citations ================
\usepackage{fancyvrb}
\usepackage{framed}
\definecolor{shadecolor}{named}{light}

% Encadrés pour les blocs > (quotes)
\usepackage{mdframed}
\newmdenv[
  linecolor=primary,
  linewidth=2pt,
  topline=false, bottomline=false, rightline=false,
  innerleftmargin=10pt, innerrightmargin=8pt,
  innertopmargin=6pt, innerbottommargin=6pt,
  backgroundcolor=codebg,
  skipabove=8pt, skipbelow=8pt
]{quotebox}
\renewenvironment{quote}{\begin{quotebox}}{\end{quotebox}}

% ================ Titres ================
\usepackage{titlesec}
\setcounter{secnumdepth}{0}
\titleformat{\section}{\Large\bfseries\color{primary}}{\thesection.}{0.6em}{}
\titleformat{\subsection}{\large\bfseries\color{accent}}{\thesubsection}{0.6em}{}
\titleformat{\subsubsection}{\normalsize\bfseries}{\thesubsubsection}{0.6em}{}

% ================ En-têtes et pieds de page ================
\usepackage{fancyhdr}
\pagestyle{fancy}
\fancyhf{}
\fancyhead[L]{\small\color{gray}ESPRIT School of Business --- L3 Business Computing BI --- Projet Intégré}
\fancyfoot[C]{\small\thepage}
\renewcommand{\headrulewidth}{0.3pt}
\renewcommand{\footrulewidth}{0pt}

% ================ Pandoc helpers ================
\providecommand{\tightlist}{\setlength{\itemsep}{0pt}\setlength{\parskip}{0pt}}
\providecommand{\passthrough}[1]{#1}

\usepackage{calc}
\providecommand{\real}[1]{#1}


% ================ Métadonnées ================
\title{\color{primary}\textbf{Description des Données}}
\author{Aymen Ben Brik et Hela Hammami}
\date{2026}

\begin{document}

\begin{titlepage}
\centering
\vspace*{1.5cm}

{\Large\bfseries ESPRIT School of Business}\\[0.3cm]
{\large Licence 3 --- Business Computing BI}\\[2cm]

{\color{gray}\large Projet Intégré}\\[0.3cm]
{\color{gray}\large Analyse et Prédiction du Churn Client}\\[2.2cm]

\rule{0.85\linewidth}{1.2pt}\\[0.6cm]
{\Huge\bfseries\color{primary} Description des Données\par}
\rule{0.85\linewidth}{1.2pt}\\[2cm]

{\large Tuteurs du projet}\\[0.2cm]
{\Large\bfseries Aymen Ben Brik et Hela Hammami}\\[1cm]
{\large 2026}

\vfill
{\color{gray}\small Document à destination des étudiants}
\end{titlepage}

\tableofcontents
\newpage

Les données mises à votre disposition sont issues d'un système
d'information bancaire. Elles ont été \textbf{anonymisées et
généralisées} : identifiants aléatoires, libellés neutres,
regroupements, dates et montants arrondis. Elles conservent leur
cohérence d'ensemble, mais ne décrivent plus aucune personne ni aucun
site de façon exploitable.

\begin{quote}
\textbf{À retenir} : les dates ont été volontairement \textbf{décalées}
et les montants \textbf{arrondis et plafonnés}. Ils ne doivent pas être
rapprochés d'un calendrier ou d'une source externe. Raisonnez sur des
durées, des évolutions et des ordres de grandeur.
\end{quote}

\section{1. Vue d'ensemble}\label{vue-densemble}

Le jeu de données est composé de :

\begin{itemize}
\tightlist
\item
  \textbf{1 fichier principal} \texttt{data\_churn.csv}
  (\textasciitilde75 Mo, 483 536 lignes, 32 colonnes) ;
\item
  \textbf{5 fichiers de référence} au format Excel (\texttt{.xlsx}) qui
  donnent le libellé associé aux codes du fichier principal.
\end{itemize}

Il décrit 363 024 clients et 410 309 comptes. Une ligne correspond à un
\textbf{rapprochement client / compte / produit} : un client ou un
compte peut donc apparaître plusieurs fois, et environ 9 \% des lignes
correspondent à des clients sans compte renseigné. Définissez clairement
votre grain avant tout calcul.

Pour les plus gros clients, le nombre de lignes a été plafonné à 100.
\textbf{Les totaux ne représentent donc pas un encours réel} :
utilisez-les pour comparer, pas comme des montants absolus.

\section{2. Fichier principal --- Description des
colonnes}\label{fichier-principal-description-des-colonnes}

\subsection{2.1 Identifiants}\label{identifiants}

{\def\LTcaptype{none} % do not increment counter
\begin{longtable}[]{@{}
  >{\raggedright\arraybackslash}p{(\linewidth - 4\tabcolsep) * \real{0.2941}}
  >{\raggedright\arraybackslash}p{(\linewidth - 4\tabcolsep) * \real{0.1765}}
  >{\raggedright\arraybackslash}p{(\linewidth - 4\tabcolsep) * \real{0.5294}}@{}}
\toprule\noalign{}
\begin{minipage}[b]{\linewidth}\raggedright
Colonne
\end{minipage} & \begin{minipage}[b]{\linewidth}\raggedright
Type
\end{minipage} & \begin{minipage}[b]{\linewidth}\raggedright
Description
\end{minipage} \\
\midrule\noalign{}
\endhead
\bottomrule\noalign{}
\endlastfoot
\texttt{CUSTOMER\_NO} & Code (\texttt{K…}) & Identifiant aléatoire du
client \\
\texttt{ACCOUNT\_NO} & Code (\texttt{X…}) & Identifiant aléatoire du
compte \\
\texttt{BRANCH\_GROUP} & Code (\texttt{AG01} à \texttt{AG20}) &
Regroupement anonyme d'agences, sans signification géographique.
\texttt{AG00} = masqué \\
\end{longtable}
}

\subsection{2.2 Informations client}\label{informations-client}

{\def\LTcaptype{none} % do not increment counter
\begin{longtable}[]{@{}
  >{\raggedright\arraybackslash}p{(\linewidth - 4\tabcolsep) * \real{0.2941}}
  >{\raggedright\arraybackslash}p{(\linewidth - 4\tabcolsep) * \real{0.1765}}
  >{\raggedright\arraybackslash}p{(\linewidth - 4\tabcolsep) * \real{0.5294}}@{}}
\toprule\noalign{}
\begin{minipage}[b]{\linewidth}\raggedright
Colonne
\end{minipage} & \begin{minipage}[b]{\linewidth}\raggedright
Type
\end{minipage} & \begin{minipage}[b]{\linewidth}\raggedright
Description
\end{minipage} \\
\midrule\noalign{}
\endhead
\bottomrule\noalign{}
\endlastfoot
\texttt{NATIONALITY\_REGION} & Texte & Zone de nationalité :
\texttt{LOCAL}, \texttt{EUROPE}, \texttt{AFRIQUE\_MO},
\texttt{AMERIQUES}, \texttt{ASIE\_OCEANIE}, \texttt{AUTRE} \\
\texttt{RESIDENCE\_REGION} & Texte & Zone de résidence (mêmes
modalités) \\
\texttt{MARITAL\_STATUS} & Code & Statut marital (C = célibataire, M =
marié(e), D = divorcé(e), V = veuf/veuve) \\
\texttt{BIRTH\_PERIOD} & Année & Début de la tranche de naissance de 5
ans (\texttt{1980} = 1980--1984). Dernière tranche basse : \texttt{1935}
et avant \\
\texttt{CUST\_OPENING\_DATE} & Date (YYYYMMDD) & Entrée en relation :
début de trimestre, d'année ou de période de 5 ans selon le client \\
\texttt{NATURE\_CLIENT} & Code & Type de client : \texttt{PPH} (personne
physique), \texttt{PM} (personne morale), \texttt{PRO} (professionnel),
\texttt{TRPP} et \texttt{TRPM} (autres catégories) \\
\texttt{PARTYCLASS} & Texte & Classe commerciale : \texttt{Retail},
\texttt{Premium}, \texttt{Corporate\ Small}, \texttt{Corporate} \\
\texttt{LOB} & Code & Ligne métier (FK vers \texttt{TARGET.xlsx}) \\
\texttt{SCORE\_KYC} & Code & Score de connaissance client (Know Your
Customer) \\
\texttt{COMPLETED\_FILE} & YES / vide & Dossier client complet \\
\texttt{LAST\_REVIEW\_DATE} & Date (YYYYMMDD) & Dernière revue du
dossier, ramenée au mois \\
\texttt{NEXT\_\_REVIEW\_DATE} & Date (YYYYMMDD) & Prochaine revue
prévue, ramenée au mois \\
\texttt{INDUSTRY} & Code (A à Z) & Section d'activité (FK vers
\texttt{INDUSTRY.xlsx}) \\
\texttt{SALARY} & Numérique & Revenu déclaré, arrondi à 1 chiffre
significatif et plafonné \\
\end{longtable}
}

\subsection{2.3 Informations compte}\label{informations-compte}

{\def\LTcaptype{none} % do not increment counter
\begin{longtable}[]{@{}
  >{\raggedright\arraybackslash}p{(\linewidth - 4\tabcolsep) * \real{0.2941}}
  >{\raggedright\arraybackslash}p{(\linewidth - 4\tabcolsep) * \real{0.1765}}
  >{\raggedright\arraybackslash}p{(\linewidth - 4\tabcolsep) * \real{0.5294}}@{}}
\toprule\noalign{}
\begin{minipage}[b]{\linewidth}\raggedright
Colonne
\end{minipage} & \begin{minipage}[b]{\linewidth}\raggedright
Type
\end{minipage} & \begin{minipage}[b]{\linewidth}\raggedright
Description
\end{minipage} \\
\midrule\noalign{}
\endhead
\bottomrule\noalign{}
\endlastfoot
\texttt{ACCOUNT\_STATUS} & Texte & \texttt{Active} ou \texttt{Closed}
--- \textbf{base de la variable cible} \\
\texttt{ACCT\_OPENING\_DATE} & Date (YYYYMMDD) & Ouverture du compte,
ramenée au lundi de la semaine \\
\texttt{ACCT\_CLOSE\_DATE} & Date (YYYYMMDD) & Clôture du compte,
ramenée au lundi de la semaine \\
\texttt{ACCOUNT\_CATEGORY} & Code (\texttt{CAT…}) & Famille de comptes
(FK vers \texttt{ACCOUNT\_CATEGORY.xlsx}) \\
\texttt{ACCOUNT\_TYPE\_DESC} & Texte & Libellé de la famille de
comptes \\
\texttt{CURRENCY} & Code & Devise (FK vers \texttt{CURRENCY.xlsx}) \\
\texttt{CLOSURE\_REASON} & Code (\texttt{RSN…}) & Motif de clôture (FK
vers \texttt{CLOSURE\_REASON.xlsx}) \\
\texttt{ACCT\_BALANCE} & Numérique & Solde arrondi à 2 chiffres
significatifs et plafonné \\
\end{longtable}
}

\subsection{2.4 Informations produit}\label{informations-produit}

{\def\LTcaptype{none} % do not increment counter
\begin{longtable}[]{@{}
  >{\raggedright\arraybackslash}p{(\linewidth - 4\tabcolsep) * \real{0.2941}}
  >{\raggedright\arraybackslash}p{(\linewidth - 4\tabcolsep) * \real{0.1765}}
  >{\raggedright\arraybackslash}p{(\linewidth - 4\tabcolsep) * \real{0.5294}}@{}}
\toprule\noalign{}
\begin{minipage}[b]{\linewidth}\raggedright
Colonne
\end{minipage} & \begin{minipage}[b]{\linewidth}\raggedright
Type
\end{minipage} & \begin{minipage}[b]{\linewidth}\raggedright
Description
\end{minipage} \\
\midrule\noalign{}
\endhead
\bottomrule\noalign{}
\endlastfoot
\texttt{PRODUCT\_GROUP} & Texte & Groupe de produits générique \\
\texttt{PRODUCT\_LINE} & Texte & Ligne de produits \\
\texttt{STARTDATE} & Date (YYYYMMDD) & Début du produit, ramené au
mois \\
\texttt{MATURITYDATE} & Date (YYYYMMDD) & Échéance, ramenée au mois \\
\texttt{AMOUNT} & Numérique & Montant du produit, arrondi à 2 chiffres
significatifs et plafonné \\
\texttt{FIXEDRATE} & Numérique & Taux fixe, arrondi au demi-point \\
\texttt{PRODUCT\_STATUS} & Texte & Statut du produit \\
\end{longtable}
}

\section{3. Fichiers de référence}\label{fichiers-de-ruxe9fuxe9rence}

{\def\LTcaptype{none} % do not increment counter
\begin{longtable}[]{@{}
  >{\raggedright\arraybackslash}p{(\linewidth - 4\tabcolsep) * \real{0.3750}}
  >{\raggedright\arraybackslash}p{(\linewidth - 4\tabcolsep) * \real{0.1875}}
  >{\raggedright\arraybackslash}p{(\linewidth - 4\tabcolsep) * \real{0.4375}}@{}}
\toprule\noalign{}
\begin{minipage}[b]{\linewidth}\raggedright
Fichier
\end{minipage} & \begin{minipage}[b]{\linewidth}\raggedright
Clé
\end{minipage} & \begin{minipage}[b]{\linewidth}\raggedright
Contenu
\end{minipage} \\
\midrule\noalign{}
\endhead
\bottomrule\noalign{}
\endlastfoot
\texttt{ACCOUNT\_CATEGORY.xlsx} & \texttt{CATEGORY\_ID} & Familles de
comptes \\
\texttt{CLOSURE\_REASON.xlsx} & \texttt{REASON\_ID} & Motifs de
clôture \\
\texttt{CURRENCY.xlsx} & \texttt{CURRENCY\_CODE} & Devises (\texttt{LCY}
= monnaie locale, \texttt{LCC} = monnaie locale convertible) \\
\texttt{INDUSTRY.xlsx} & \texttt{INDUSTRY\_CODE} & Sections
d'activité \\
\texttt{TARGET.xlsx} & \texttt{TARGET\_CODE} & Lignes métier \\
\end{longtable}
}

Le calendrier n'est pas fourni : il est à construire vous-même.

\section{4. Définition de la variable cible
(churn)}\label{duxe9finition-de-la-variable-cible-churn}

Le \textbf{churn} se définit ici comme la \textbf{clôture d'un compte} :

\begin{verbatim}
churn = 1  si  ACCOUNT_STATUS = "Closed"
churn = 0  si  ACCOUNT_STATUS = "Active"
ligne exclue si ACCOUNT_STATUS est vide
\end{verbatim}

Une définition plus fine est possible (par exemple : churn d'un client =
clôture de \textbf{tous} ses comptes), à justifier dans le rapport. Le
taux dépend fortement du grain : environ \textbf{40 \%} des lignes,
\textbf{36 \%} des comptes et \textbf{32 \%} des clients. Le
déséquilibre reste donc modéré : vérifiez-le et justifiez toute
technique de rééquilibrage.

\section{5. Variables pour la
prédiction}\label{variables-pour-la-pruxe9diction}

\subsection{5.1 Attention à la fuite de
cible}\label{attention-uxe0-la-fuite-de-cible}

Le fichier est une \textbf{photographie} prise après les clôtures :
certaines colonnes révèlent le résultat au lieu de le prédire. \textbf{À
ne pas utiliser comme variables explicatives} :

\begin{itemize}
\tightlist
\item
  \texttt{ACCT\_CLOSE\_DATE} et \texttt{CLOSURE\_REASON} : renseignées
  uniquement pour les comptes clos.
\item
  \texttt{ACCT\_BALANCE} \textbf{et sa valeur manquante} : le solde est
  vide pour environ 68 \% des comptes clos et n'est jamais vide pour les
  comptes actifs. Toute variable dérivée de ce vide trahit la cible.
\item
  La famille de compte, la devise et le montant : leur disponibilité
  peut suivre le même schéma, à contrôler.
\item
  \texttt{PRODUCT\_STATUS} : certaines modalités correspondent à la
  fermeture.
\end{itemize}

Un score proche de 100 \% signale presque toujours une fuite, pas un bon
modèle. Le rapport doit présenter ce contrôle.

\subsection{5.2 Pistes de variables}\label{pistes-de-variables}

\begin{itemize}
\tightlist
\item
  \textbf{Démographiques} : tranche d'âge, statut marital, zone de
  nationalité ou de résidence.
\item
  \textbf{Ancienneté} : ancienneté du client et du compte, à calculer
  depuis leurs dates d'ouverture.
\item
  \textbf{Comportementales} : nombre de comptes par client, diversité
  des produits, secteur, segment commercial.
\item
  \textbf{Conformité} : score KYC, dossier complet, ancienneté de la
  dernière revue.
\end{itemize}

\section{6. Séries temporelles --- points
d'attention}\label{suxe9ries-temporelles-points-dattention}

\texttt{ACCT\_CLOSE\_DATE} et \texttt{ACCT\_OPENING\_DATE} permettent de
construire des séries \textbf{hebdomadaires}. La granularité mensuelle
ne suffit pas : la période de clôtures couvre environ 47 mois, soit trop
peu de points pour un LSTM.

\begin{itemize}
\tightlist
\item
  La série des clôtures compte environ \textbf{200 semaines}, dont une
  dizaine sans aucune clôture : ces semaines valent \textbf{0} et
  doivent être ajoutées.
\item
  \textbf{Cinq semaines concentrent environ la moitié des clôtures} : ce
  sont des opérations de masse. Elles doivent être identifiées et
  traitées (variable d'intervention, série corrigée, modèle robuste),
  puis ce choix doit être justifié.
\item
  Comptez la série au grain \textbf{compte} après déduplication, sinon
  les lignes multiples la déforment.
\item
  Le découpage entraînement / test doit \textbf{respecter l'ordre du
  temps}. Aucun mélange aléatoire.
\end{itemize}

\section{7. Qualité des données}\label{qualituxe9-des-donnuxe9es}

Le jeu conserve des défauts à traiter dans vos phases d'analyse et d'ETL
:

\begin{itemize}
\tightlist
\item
  \textbf{Valeurs manquantes} nombreuses : plus de 85 \% sur
  \texttt{MATURITYDATE}, \texttt{FIXEDRATE}, \texttt{STARTDATE},
  \texttt{PRODUCT\_STATUS} ; environ 70 \% sur \texttt{SALARY}. Certains
  vides sont \textbf{volontaires} et viennent de l'anonymisation.
\item
  \textbf{Doublons} : plusieurs milliers de lignes strictement
  identiques, et bien plus au grain compte.
\item
  \textbf{Dates numériques} au format \texttt{YYYYMMDD} à convertir. La
  précision de \texttt{CUST\_OPENING\_DATE} varie selon les clients.
\item
  \textbf{Soldes négatifs} et \textbf{montants extrêmes} : environ 32 \%
  des soldes sortent des bornes IQR. Les comptes débiteurs expliquent
  une partie de ces valeurs, qu'il faut interpréter avant de les
  exclure.
\item
  \textbf{Valeurs plafonnées} : un pic sur la valeur maximale d'un
  montant vient du plafonnement, pas d'un phénomène métier.
\end{itemize}

\section{8. Précautions}\label{pruxe9cautions}

\begin{itemize}
\tightlist
\item
  \textbf{Ne pas publier les données}, sur GitHub ou tout autre service.
\item
  \textbf{Ne pas chercher} à identifier l'établissement, une agence ou
  un client.
\item
  Utiliser un \texttt{.gitignore} qui exclut au minimum :
  \texttt{data/}, \texttt{*.csv}, \texttt{*.xlsx}, \texttt{*.pkl},
  \texttt{*.pbix}.
\item
  Ne pas envoyer les données à un service externe (assistant en ligne,
  outil de profilage hébergé) sans l'accord des tuteurs.
\end{itemize}

\end{document}
